\section{Giới thiệu}
\label{sec:introduction}

Các mô hình ngôn ngữ lớn (Large Language Models -- LLMs) đã đạt nhiều kết quả
tích cực trong hỏi đáp. Tuy nhiên, trong lĩnh vực pháp luật, một câu trả lời
đáng tin cậy không chỉ cần trôi chảy mà còn phải xác định đúng căn cứ và chuỗi
suy luận nối các điều kiện, sự kiện trung gian và hệ quả
\cite{chalkidis2020legalbert,guha2023legalbench}.

Retrieval-Augmented Generation (RAG) giảm phụ thuộc vào tri thức nội tại của
LLM bằng cách truy xuất bằng chứng trước bước sinh
\cite{lewis2020retrieval,gao2024ragsurvey}. Tuy nhiên, RAG truyền thống chủ yếu
dựa vào tương đồng ngữ nghĩa và thường xem các đoạn văn là bằng chứng độc lập.
Cách này còn hạn chế với câu hỏi pháp luật nhiều bước, trong đó kết luận phụ
thuộc vào chuỗi quy tắc điều kiện--hệ quả liên kết.

Chúng tôi đề xuất \textbf{LegalSCM-RAG}, kết hợp truy xuất trên đồ thị pháp luật
có hướng với xác minh cấu trúc phụ thuộc câu hỏi. Verifier chuyển claim trực
tiếp tới phép kiểm tra cạnh đã ground, câu hỏi thông thường tới suy luận factual
ba giá trị, và câu hỏi nêu rõ việc loại mediator tới hard intervention. Commit
gate chỉ sử dụng quyết định verifier khi thỏa mãn các yêu cầu về grounding, độ
phủ đường, confidence và nhất quán trạng thái; nếu không, hệ thống giữ nguyên
context và quyết định của Causal Path RAG.

Các đóng góp chính gồm:
\begin{itemize}
    \item formulation điểm bất động ba giá trị với mã hóa incidence thưa cho
    thân/đầu quy tắc, trong đó conjunction nằm trong một quy tắc và disjunction
    thay thế nằm giữa nhiều quy tắc;
    \item phân biệt chính thức và ở mức implementation giữa node-deletion
    reachability với mechanism-replacement $do(X_M=0)$; và
    \item đánh giá có kiểm soát bốn hệ thống, kèm kiểm định theo cặp, phân tích
    abstention và audit route/commit.
\end{itemize}

Trên benchmark hiện tại, selective verification cải thiện có ý nghĩa Decision
Accuracy so với Causal Path RAG nhưng không cải thiện chất lượng answer hoặc
citation. Batch được báo cáo chỉ gồm factual validation và direct-edge route,
không có mediator intervention nào được thực thi. Vì vậy, mức tăng được diễn
giải là hiệu quả của toàn bộ selective verifier, không phải bằng chứng thực
nghiệm riêng cho structural counterfactual intervention.
